\begin{claim}{C.7e}{Proved}
Let $K/\Q_p$ be finite, let $k_0\subset K$ be a subfield finitely
generated over $\Q$, and let $U_0/k_0$ be a hyperbolic curve. Put
$U=U_0\times_{k_0}K$, $\Pi=\pi_1(U_{\bar K})$, and
$H=G_{K\Q_p(\mu_\infty)}$.

If an ordinary $x\in U(K)$ is not algebraic over $k_0$---equivalently,
its image in the curve $U_0$ is the generic point---then
\[
                         \Pi^{s_x(H)}=1.                     \tag{C.7e.1}
\]
If $c\in(\bar U_0\setminus U_0)(k_0)$, $u$ is a local parameter
at $c$ defined over $k_0$, and the tangent is specified by
$u=a t$ with $a\in K^\times$ transcendental over $k_0$, then
\[
                  \Pi^{s_{c,a}(H)}=I_c.                    \tag{C.7e.2}
\]
Both equalities remain true after replacing $H$ by any open subgroup.
These are equalities in the full profinite group.

For example, on the tripod over $\Q_5$ let
\[
 \theta=\sum_{n=1}^{\infty}5^{\,n!},\qquad x=2+\theta,
 \qquad a=1+\theta.
\]
The point $x$ specializes to the ordinary point $2$, its full
cyclotomic fixed group is trivial, and the integral tangent with
direction coefficient $a$ at $0$ has full fixed group $I_0$.
The theorem does not decide the algebraic point $2$ or the tangent
with coefficient $1$.
\end{claim}
\begin{proof}
\emph{A local abelian restriction lemma.}
For every finite $M/\Q_p$ and every abelian variety $A/M$, one has
\[
 H^1(G_M,V_pA)\ \hookrightarrow\
 H^1(G_{M\Q_p(\mu_\infty)},V_pA).                            \tag{C.7e.3}
\]
This assertion does not require good reduction of $A$.
Write $H_M=G_{M\Q_p(\mu_\infty)}$,
$\Gamma_M=G_M/H_M$, and $W=(V_pA)^{H_M}$.
The quotient $\Gamma_M$ is abelian. Mattuck's theorem gives an open
subgroup of finite index in $A(M)$ isomorphic to the additive group
$\mathcal O_M^{\dim A}$
\cite[Theorem~7, printed pp.~114--117]{BC_Mattuck1955}.
In particular $A(M)_{\rm tors}$ is finite, so
\[
                        W^{\Gamma_M}=(V_pA)^{G_M}=0.          \tag{C.7e.4}
\]
Indeed a nonzero fixed vector could be scaled into $T_pA$ and
would give $M$-rational torsion of unbounded $p$-power order.

Here is the continuous-cohomology argument for
$H^1(\Gamma_M,W)=0$, with no semisimplicity assumption.
Let $\mathcal B$ be the finite-dimensional commutative
$\Q_p$-algebra generated by the action of $\Gamma_M$ on $W$.
After a finite coefficient extension $E/\Q_p$, its semisimple
quotient splits and
\[
                         W_E=\bigoplus_\lambda W_\lambda
\]
is its decomposition into joint generalized-character spaces.
On $W_\lambda$, every group element has scalar semisimple
part $\lambda(\gamma)$. A block with $\lambda=1$ would contain
a common fixed vector: its local-algebra radical is nilpotent,
and a last nonzero power of that radical applied to $W_\lambda$
is annihilated by every $\gamma-1$. This contradicts (C.7e.4).
Thus for every block there is a $\gamma_\lambda\in\Gamma_M$
such that $\gamma_\lambda-1$ is invertible on $W_\lambda$.

For a continuous $1$-cocycle $b:\Gamma_M\to W_\lambda$,
commutativity gives
\[
  (\gamma_\lambda-1)b(\sigma)
                   =(\sigma-1)b(\gamma_\lambda).
\]
Consequently $b(\sigma)=(\sigma-1)v_\lambda$, where
$v_\lambda=(\gamma_\lambda-1)^{-1}b(\gamma_\lambda)$.
Taking the direct sum shows that every cocycle becomes a
coboundary over $E$. Normalized trace $E\to\Q_p$ descends its
cobounding vector, so it was already a coboundary over $\Q_p$.

To deduce (C.7e.3) directly, take a cocycle on $G_M$ whose
restriction to $H_M$ is a coboundary, and subtract that
coboundary. The resulting cocycle is zero on $H_M$.
The cocycle identity shows that its values lie in $W$ and that
it descends continuously to $\Gamma_M$. The vanishing just
proved kills it.

The kernel of the rational $p$-adic Kummer map
\[
                    A(M)\longrightarrow H^1(G_M,V_pA)
\]
is exactly $A(M)_{\rm tors}$. One can see this directly from
Mattuck's finite-index subgroup. The compact abelian group
$A(M)$ is a finite extension of a finite-rank free $\Z_p$-module;
its $p$-adic completion, tensored with $\Q_p$, therefore kills
exactly its torsion. The integral Kummer map on that completion
is injective: at level $p^n$ its kernel is $p^n A(M)$, by the
multiplication-by-$p^n$ exact sequence. Alternatively, if an
integral Kummer cocycle becomes a coboundary with values in
$V_pA$, multiply its cobounding vector by a power of $p$ to
put it in $T_pA$. That multiple of the point then lies in
$\bigcap_n p^n A(M)$, which is finite torsion by the same
finite-index description.

It follows that an $H_M$-invariant path between ordinary points
$P,Q$ of any smooth proper curve over $M$ forces
\[
                       [P-Q]\in\operatorname{Jac}(C)(M)
                                \quad\hbox{to be torsion}.    \tag{C.7e.5}
\]
Push the path through the Abel map based at $P$ and apply
(C.7e.3). The resulting path torsor is the Kummer torsor of
$[Q-P]$. The same implication applies to tangential sections
after filling in their underlying cusps. We use only torsion
of this particular point over the finite local field $M$;
no finiteness theorem for $A(M(\mu_\infty))_{\rm tors}$ is
assumed.

\emph{Descent of the geometric covering data.}
Let $L$ be the relative algebraic closure of $k_0$ in $\bar K$.
It is algebraically closed. For the smooth curve $U_{0,L}$,
the extension $L\subset\bar K$ induces an equivalence of
finite etale covering categories. This is the characteristic-zero
case of SGA~1, XIII, Proposition~4.6, taking the second factor
there to be $\operatorname{Spec}\bar K$
\cite[XIII, Proposition~4.6, original pp.~421--422]{BC_SGA1}.
The resolution hypothesis in that statement is needed only in
dimension at most one here. All primes are retained in
characteristic zero.

Thus any finite geometric cover used below, its Galois closure,
its compactification, and its finitely many deck transformations
can be defined over $L$. Each finite collection of such data
descends further to a finite extension of $k_0$, and hence can
be defined over a finite local extension of $K$. For an abelian
variety defined over $L$, every torsion point over $\bar K$
already lies over $L$: its order-$n$ subgroup is a finite
$L$-scheme and $L$ is algebraically closed.

\emph{A refinement with proper genus at least two.}
Every hyperbolic curve over $L$ has a connected finite etale
cover whose compactification has genus at least two.
If the original proper genus is at least two, use the identity.
In genus zero, retain three cusps and use the Fermat-quartic
cover
\[
        r^4+s^4=1,\qquad x=r^4
\]
of the tripod. It is a connected Galois cover of degree $16$,
ramified only over $0,1,\infty$, with ramification index four
at each. Its compactification has genus three by
$2g-2=-32+3\cdot4\cdot3=4$. Restrict it to the given
punctured projective line.

In genus one there is at least one cusp. Choose an origin on
the compact elliptic curve and first pull back by multiplication
by two. The resulting punctured elliptic curve has at least
four distinct cusps. Choose two of them, $A,B$, and choose
$R$ with $2R=A+B$ in the elliptic group. A function with
divisor $A+B-2R$ exists. Its square-root cover is connected,
ramified only at $A,B$, and has compact genus two. It is
etale over the punctured curve. These constructions all take
place in characteristic zero. Taking a Galois closure and then
a compositum with any prescribed finite Galois cover gives a
Galois refinement of that cover with proper genus at least two:
a finite map to a genus-at-least-two curve has source genus
at least two by Riemann--Hurwitz.

\emph{Ordinary points not algebraic over the field of definition.}
Suppose $g\in\Pi^{s_x(H)}$ is nontrivial. Choose a finite
Galois quotient detecting $g$ and refine it as in the preceding
paragraph. Let $Y\to U$ be the resulting Galois cover, let
$D$ be its compactification, and let $G$ be its deck group;
then $g(D)\ge2$ and the image of $g$ in $G$ is nontrivial.
Choose a point $P$ of $Y_{\bar K}$ over $x$ and lift the loop
$g$ from $P$ to $Q$. The fibre is a free $G$-set, so
$Q=h(P)$ for a nonidentity $h\in G$. Choose a finite local
extension $M/K$ defining all selected covering data and endpoints.
The lifted path is $H_M$-invariant by uniqueness of path lifting.
Formula (C.7e.5) gives
\[
                              [P-h(P)]\ \text{torsion}.
\]
The curve $D$, the automorphism $h$, and its Jacobian are defined
over $L$, so this torsion value belongs to $\operatorname{Jac}(D)(L)$.

The $L$-morphism
\[
        d_h:D\longrightarrow\operatorname{Jac}(D),
                    \qquad R\longmapsto[R-h(R)]
\]
is nonconstant. Indeed, if it were constant, then for a fixed
$O\in D(L)$,
\[
     [h(R)-h(O)]=[R-O]\qquad\hbox{for every }R,
\]
so $h$ would act trivially on the Jacobian, since the Abel
image generates it. An automorphism of a smooth proper
characteristic-zero curve of genus at least two acting trivially
on its Jacobian is the identity. To verify this last assertion,
it acts trivially on all regular differentials. In the
nonhyperelliptic case the canonical map is an embedding.
In the hyperelliptic case it would be either the identity or
the hyperelliptic involution, and the latter acts by $-1$ on
the regular differentials. Thus neither case allows this $h$.

A nonconstant morphism from a proper integral curve has finite
fibres. Its fibre over the selected $L$-rational torsion value
is therefore a finite $L$-scheme, so every geometric point
of that fibre, in particular $P$, is $L$-rational. Its image
$x$ is then algebraic over $k_0$, a contradiction. This
proves (C.7e.1).

\emph{Transcendental tangent directions.}
The inertia $I_c$ is fixed by $s_{c,a}(H)$, because its
conjugation action is cyclotomic. Suppose that a full fixed
element $g$ does not belong to $I_c$. Since $I_c$ is closed,
there is an open normal subgroup $N\triangleleft\Pi$ with
$g\notin N I_c$. Let $Y\to U$ be the corresponding Galois
cover and $D$ its compactification, all defined over $L$.
The lift of $g$ from a chosen tangent ends at another tangent
whose underlying cusp is distinct: in the regular fibre,
cusps above $c$ are the right orbits of the image of $I_c$,
and the image of $g$ lies outside that subgroup.
Call the two underlying cusps $P,Q$.

Extend $K$ finitely to a local $M$ containing all coefficients
and the chosen lifted tangents. The path is then $H_M$-invariant.
After filling the boundary, (C.7e.5) makes $[P-Q]$ torsion.
Thus for some integer $m>0$ there is a rational function
$f\in L(D)^\times$ with
\[
                           \operatorname{div}(f)=m(P-Q).
                                                               \tag{C.7e.6}
\]
For proper genus zero this conclusion holds directly with $m=1$.
In all cases it can be taken over $L$, by the $L$-rational
Riemann--Roch space of the divisor. Its restriction to $Y$
is an invertible function.

Choose local parameters $z,w$ at $P,Q$ over $L$. Since the
cover is Galois, the two ramification indices above $c$ are
the same integer $e>0$. Write
\[
 \begin{split}
       u&=c_P z^e(1+O(z))=c_Q w^e(1+O(w)),\\
       f&=d_P z^m(1+O(z))\quad\hbox{at }P,\\
       f&=d_Q w^{-m}(1+O(w))\quad\hbox{at }Q,
 \end{split}                                                  \tag{C.7e.7}
\]
with $c_P,c_Q,d_P,d_Q\in L^\times$.
The tangent $u=a t$ lifts with leading coefficients
\[
          z=\alpha\,t^{1/e}+\cdots,\qquad
          w=\beta\,t^{1/e}+\cdots,\qquad
          \alpha^e=a/c_P,\quad \beta^e=a/c_Q.                 \tag{C.7e.8}
\]
All selected constants and $\alpha,\beta$ can be included
in the finite local field $M$. The roots of the higher
unit terms in the expansions are uniquely obtained by
formal power series in characteristic zero. Thus the
Kummer cocycles of the two induced fibre functors on
$\mathbf G_m$ are those of their leading coefficients
$d_P\alpha^m$ and $d_Q\beta^{-m}$. More explicitly, taking
an $n$th root of $f$ in the first Puiseux expansion gives
the constant factor $(d_P\alpha^m)^{1/n}$ times
$t^{m/(en)}$ and a uniquely normalized unit; the second
expansion gives $(d_Q\beta^{-m})^{1/n}$ times
$t^{-m/(en)}$. This proves the asserted coefficient
comparison without identifying the two boundary fibre
functors with ordinary points of $\mathbf G_m$.

The $H_M$-invariant path pushed through the unit $f$
therefore forces
\[
       \kappa_p\!\left((d_P/d_Q)(\alpha\beta)^m\right)
                                      \bigm|_{H_M}=0.
                                                               \tag{C.7e.9}
\]
The multiplicative restriction lemma (C.7b.1) makes the
displayed ratio a root of unity. Such a root belongs to
$L$, so $(\alpha\beta)^m\in L^\times$.
Raising to the $e$th power and using (C.7e.8) gives
\[
                  a^{2m}=(c_Pc_Q)^m(\alpha\beta)^{em}
                                      \in L^\times.
\]
Since $L$ is algebraically closed in $\bar K$, this
forces $a\in L$, contrary to the transcendence assumption.
Hence every fixed element lies in $I_c$, proving (C.7e.2).

Every open subgroup of $H$ is $H_M$ for some finite local
extension $M/K$, as in C.7d. The hypotheses of the two
arguments are unchanged by that extension, proving the
assertion about open subgroups.

\emph{Explicit parameters.}
The series $\theta$ converges in $5\Z_5$. It is
transcendental over $\Q$. If a nonzero integer polynomial
$F$ of degree $d$ vanished at $\theta$, put
$r_N=\sum_{n=1}^N5^{n!}$. For all sufficiently large $N$,
$F(r_N)\ne0$, and the $5$-adic polynomial difference
identity would give
\[
               v_5(F(r_N))\ge (N+1)!.
\]
On the other hand $0<r_N<5^{N!+1}$, so for a fixed
constant $C$,
\[
  5^{(N+1)!}\le |F(r_N)|_{\mathbf R}
                         \le C\,5^{d(N!+1)},
\]
which is impossible as $N\to\infty$.
Thus both $x=2+\theta$ and $a=1+\theta$ meet the theorem's
transcendence hypotheses. Their reductions modulo five
are respectively $2$ and $1$, as asserted.
\end{proof}
